% ==============================================================================
% Chapitre 6 : Extensions et Évolutions Futures
% Roadmap post-projet, évolutions potentielles, améliorations continues
% ==============================================================================

\chapter{Extensions et Évolutions Futures}
\label{ch:extensions}

\section{Roadmap post-projet}

Le projet LockBits a livré une version fonctionnelle complète (v1.0.0) qui
couvre les besoins essentiels~: serveur EDR avec analyse de fichiers,
agent déployable, portail client, pipeline CI/CD et infrastructure de
monitoring. Cette section présente la roadmap post-MVP, c'est-à-dire les
axes d'amélioration identifiés pour renforcer la robustesse, la
couverture fonctionnelle et la qualité du logiciel.

\subsection{Court terme (post-MVP)}

Les priorités à court terme concernent les correctifs et les
améliorations de qualité identifiés durant le développement de la
version initiale~:

\begin{enumerate}
  \item \textbf{Correction du bug de création de ticket GLPI}~: un
    dysfonctionnement dans l'appel API GLPI côté site client empêche la
    création de tickets de support. Ce bug bloque une fonctionnalité
    métier importante et doit être résolu en priorité.
  \item \textbf{Amélioration de la couverture de tests unitaires}~:
    l'objectif est d'atteindre 80~\% de couverture de code pour le
    serveur EDR. Une couverture élevée réduit le risque de régressions
    et facilite les évolutions futures.
  \item \textbf{Tests d'intégration pour le site client}~: des tests
    automatisés spécifiques au portail PHP doivent être ajoutés,
    couvrant les scénarios d'authentification, de consultation du
    tableau de bord et d'interaction avec l'API EDR.
  \item \textbf{Backup automatisé de la base MySQL}~: un script de
    sauvegarde avec rotation doit être mis en place pour la base de
    données du portail client. La politique de rétention prévoit la
    conservation des sept derniers snapshots quotidiens et des quatre
    dernières sauvegardes hebdomadaires.
  \item \textbf{Configuration Dokploy finalisée}~: s'assurer que la configuration
      Dokploy couvre tous les environnements (prod et preprod) de manière robuste.
\end{enumerate}

\subsection{Moyen terme}

Les améliorations à moyen terme visent à renforcer la supervision,
l'accessibilité et la fiabilité de la plateforme~:

\begin{enumerate}
  \item \textbf{Alerting Prometheus}~: la configuration des alertes
    Prometheus doit être enrichie avec des webhooks vers Discord et par
    email. Les administrateurs seront notifiés en cas de dysfonctionnement
    d'un service, de saturation d'un disque ou d'absence prolongée d'un
    agent EDR.
  \item \textbf{Support IPv6}~: la stack Docker doit supporter
    l'adressage IPv6 en complément de l'IPv4. Cette évolution est
    nécessaire pour garantir la compatibilité avec les infrastructures
    réseau modernes qui déploient progressivement IPv6.
  \item \textbf{Rate limiting avancé sur l'API EDR}~: le mécanisme de
    rate limiting actuel doit être enrichi avec des politiques
    différenciées par client, des limites ajustables dynamiquement et
    une meilleure visibilité sur les quotas consommés via les endpoints
    de monitoring.
  \item \textbf{Dashboard Grafana prêt à l'emploi}~: le dashboard de
    monitoring doit être importé automatiquement au déploiement, sans
    intervention manuelle. Cela garantit que chaque instance déployée
    dispose immédiatement d'une visibilité complète sur ses métriques.
\end{enumerate}

\section{Évolutions potentielles}

Au-delà des améliorations immédiates, plusieurs évolutions
architecturales et fonctionnelles ont été identifiées pour les phases
ultérieures du projet. Ces évolutions répondent à des besoins plus
ambitieux et nécessitent des investissements plus importants.

\subsection{Déploiement Kubernetes (K3s)}

La migration de Docker Compose vers Kubernetes représente l'évolution
la plus significative de l'infrastructure. \textbf{K3s}, distribution
Kubernetes légère de la CNCF, est particulièrement adapté aux petites
infrastructures et aux environnements aux ressources limitées.

Les avantages de cette migration sont nombreux~:

\begin{itemize}
  \item \textbf{Orchestration avancée}~: Kubernetes gère le cycle de vie
    des conteneurs, les redémarrages automatiques et la répartition de
    charge entre les nœuds.
  \item \textbf{Auto-scaling}~: le nombre de réplicas de chaque service
    peut s'ajuster automatiquement en fonction de la charge, ce qui
    permet de faire face aux pics d'activité sans sur-allouer les
    ressources.
  \item \textbf{Rolling updates}~: les mises à jour sont appliquées
    progressivement sans interruption de service. Kubernetes remplace
    les pods un par un en vérifiant la santé de chaque nouvelle
    instance.
  \item \textbf{Self-healing}~: en cas de défaillance d'un conteneur ou
    d'un nœud, Kubernetes le redémarre ou le remplace automatiquement
    pour maintenir le niveau de service défini.
\end{itemize}

Cette migration présente néanmoins des défis notables~: la complexité
opérationnelle est plus élevée, la courbe d'apprentissage pour l'équipe
est significative, et la consommation de ressources est supérieure à
celle de Docker Compose. Une phase de pré-production prolongée serait
nécessaire pour valider la configuration avant la migration en
production.

\subsection{Authentification OAuth2 complète}

Le système d'authentification actuel repose sur l'intégration GLPI avec
un mécanisme de SSO propriétaire. Une évolution vers OAuth2 standard
permettrait de supporter l'authentification via les fournisseurs
d'identité majeurs~: Google, Microsoft et GitHub.

Cette évolution apporterait plusieurs bénéfices~:

\begin{itemize}
  \item \textbf{Simplification de l'inscription}~: les nouveaux clients
    peuvent créer un compte en quelques clics via leur fournisseur
    d'identité existant, sans formulaire d'inscription.
  \item \textbf{Réduction de la friction}~: la connexion est plus rapide
    et plus sécurisée, avec des mécanismes de récupération de mot de
    passe délégués au fournisseur.
  \item \textbf{Complément à GLPI}~: OAuth2 ne remplace pas l'intégration
    GLPI mais la complète. Le SSO GLPI conserve son rôle pour la
    gestion des tickets et le support technique.
\end{itemize}

\subsection{SIEM (Wazuh / ELK)}

L'intégration d'un \textbf{SIEM} (Security Information and Event
Management) est une évolution naturelle pour un produit EDR. Le SIEM
permet de corréler les événements de sécurité remontés par les agents
avec d'autres sources de données, d'automatiser la détection de
menaces complexes et de produire des rapports de conformité.

Deux solutions open-source sont envisagées~:

\begin{itemize}
  \item \textbf{Wazuh}~: fork de l'OSSEC, Wazuh est une plateforme de
    sécurité open-source mature qui intègre la détection d'intrusion
    (HIDS), le monitoring d'intégrité de fichiers (FIM) et l'analyse
    des vulnérabilités. Son tableau de bord Kibana offre une
    visualisation complète des alertes de sécurité.
  \item \textbf{ELK Stack} (Elasticsearch, Logstash, Kibana)~: solution
    de référence pour la centralisation et l'analyse de logs.
    Elasticsearch indexe les événements, Logstash les normalise et
    Kibana permet l'exploration et la création de dashboards.
\end{itemize}

L'intégration d'un SIEM implique la définition de règles de détection
personnalisées, adaptées aux menaces spécifiques auxquelles les clients
LockBits sont exposés. Ces règles pourraient être partagées sous forme
de package de détection open-source, contribuant ainsi à l'écosystème
de la cybersécurité.

\subsection{Interface d'administration web EDR}

Le portail client actuel offre un tableau de bord orienté vers les
besoins des clients finaux. Une interface d'administration dédiée à
l'EDR est envisagée pour répondre aux besoins des équipes SOC (Security
Operations Center).

Cette interface permettrait de~:

\begin{itemize}
  \item Gérer les agents, leur déploiement et leur cycle de vie.
  \item Visualiser les alertes de sécurité avec leur niveau de gravité
    et leur statut.
  \item Générer des rapports STIX exploitables par d'autres outils de
    cybersécurité.
  \item Configurer les politiques de détection et les règles de
    réponse automatique.
  \item Suivre en temps réel l'activité des endpoints protégés.
\end{itemize}

Cette interface viendrait en complément du dashboard client existant,
offrant une séparation claire entre la vue client (statut de
l'abonnement, tickets de support) et la vue administration (opérations
de sécurité, analyses forensiques).

\subsection{Tests de pénétration automatisés}

La sécurité étant au cœur du produit LockBits, l'intégration de tests
de pénétration automatisés dans la pipeline CI/CD est une évolution
naturelle. L'objectif est de détecter les vulnérabilités avant qu'elles
n'atteignent la production.

Les axes d'automatisation envisagés sont~:

\begin{itemize}
  \item \textbf{Scans de vulnérabilités réguliers}~: au-delà de Trivy
    déjà intégré pour les images Docker, des scans de code source avec
    des outils comme Semgrep ou CodeQL permettraient de détecter les
    failles logicielles (injections SQL, XSS, failles de sécurité).
  \item \textbf{Tests d'intrusion automatisés}~: des outils comme OWASP
    ZAP peuvent être intégrés à la pipeline pour réaliser des tests
    d'intrusion sur les endpoints exposés, avec génération automatique
    de rapports.
  \item \textbf{Analyse de dépendances}~: Dependabot est déjà actif sur
    les dépôts du projet, mais une analyse plus poussée des
    dépendances avec un outil de type Snyk pourrait compléter la
    couverture de sécurité.
\end{itemize}

\section{Améliorations continues}

Au-delà des évolutions fonctionnelles et architecturales, la pérennité
du projet LockBits repose sur une démarche d'amélioration continue.
Cette section décrit les pratiques qui doivent être maintenues et
renforcées tout au long du cycle de vie du logiciel.


\subsection{Revue de code systématique}
Chaque modification du code source
doit passer par une revue par au moins un autre membre de l'équipe. La
revue vérifie la conformité aux normes de codage, la pertinence des
tests associés et l'absence d'effets de bord. Les pull requests doivent
rester de taille modérée pour faciliter la relecture.

\subsection{Mise à jour des dépendances}
Les dépendances logicielles
évoluent rapidement, et les correctifs de sécurité sont publiés
régulièrement. L'outil Dependabot, déjà configuré sur les dépôts,
propose automatiquement des mises à jour. L'équipe s'engage à les
traiter dans un délai raisonnable, en particulier pour les correctifs
de sécurité critiques.

\subsection{Documentation vivante}
La documentation technique doit évoluer
avec le logiciel. Chaque nouvelle fonctionnalité, chaque modification
d'architecture et chaque changement de configuration doit être reflété
dans les fichiers de documentation concernés. L'objectif est de
maintenir un écart minimal entre le code et sa documentation, pour
éviter que cette dernière ne devienne obsolète.

\subsection{Monitoring des performances}
\textbf{} Les métriques de performance
collectées par Prometheus et visualisées dans Grafana doivent être
revues périodiquement. Des seuils d'alerte doivent être définis pour
anticiper les dégradations avant qu'elles n'affectent les clients. Le
temps de réponse de l'API EDR, l'utilisation mémoire des conteneurs et
le nombre d'agents actifs sont les indicateurs clés à surveiller en
priorité.

\subsection{Veille technologique}
Enfin, la veille technologique fait partie intégrante de la démarche
d'amélioration continue. Le paysage de la cybersécurité évolue
rapidement~: de nouvelles menaces apparaissent, de nouveaux outils
s'imposent, et les bonnes pratiques se raffinent. L'équipe LockBits
doit rester attentive à ces évolutions pour maintenir la pertinence et
l'efficacité de sa solution face aux défis de sécurité de demain.
